\section{Discussion}\label{sec:discussion}

\subsection{Managerial implications}

\paragraph{Protect scarce clean capacity} In multi-site carbon-aware scheduling, the largest operational loss does not come from imperfect carbon forecasts. It comes from spending low-carbon capacity on work that could have waited or run elsewhere. Proposition~\ref{prop:myopic} and Theorem~\ref{thm:lb} show this in the worst case, and the experiments confirm it on average. With migration allowed, a myopic planner with perfect carbon foresight is still 17--21\% above the optimum in the synthetic experiment and 34\% above it on the trace-derived workload. Clean capacity at a low-carbon site should be treated as a scarce resource with an opportunity cost, as in revenue management. A simple protection level (MPC-Protect) captures part of this value. CARMA captures more, because its ghost jobs price the protection hour by hour from the operator's own history. The result held when the profile underestimated demand by 28\% and demand was bursty. A stochastic lookahead over sampled demand scenarios does slightly better in Great Britain and the United States (0.4--1.0 points of gap) and is level with CARMA elsewhere, at 12--37 times the decision time; when the planning network can be solved in seconds, it is a reasonable alternative. The value of demand information is not specific to one grid: the advantage of CARMA over the best baseline without it is 7.3 points in Europe and 1.7--2.8 points in the other three fleets, and it tends to be smaller where the sites are closer to each other relative to their mean intensity.

\paragraph{Invest in forecasts where they pay} The value of better intensity forecasts depends on the rest of the policy. For a myopic multi-site planner, perfect forecasts are worth only 1.4--2.4 percentage points; once demand is anticipated they are worth 4--7 points; without migration they are worth almost the whole remaining gap. Persistent errors, such as a missed day-ahead wind front, are more costly than independent ones. Forecasting effort, for example the integration of numerical weather predictions~\cite{maji2022carboncast,lowry2018day}, should therefore target day-level biases, and in multi-site settings it should follow the fix of demand myopia.

\paragraph{Use risk premia with care} The conformal cost views help a myopic planner but not CARMA. They work by making volatile sites and distant hours look dirtier, which crudely imitates anticipation. Once demand is anticipated, the same distortion gives up useful deferrals. Calibrated quantiles are better suited to reporting and to service-level decisions than to reshaping costs by default.

\paragraph{Emissions or service near saturation} Without migration and at high load, every carbon-aware policy delayed work, and under a service-adjusted metric carbon-agnostic operation was preferable. Operators should combine temporal shifting with explicit service constraints, or restrict it, when local capacity is close to saturation.

\paragraph{Attributional versus physical savings} The headline reductions are attributional: they measure how much of the work runs where and when the \emph{average} regional mix is clean. Within a single market such as Great Britain, moving load between regions mostly reassigns the average mix. Under a stylized marginal-emission proxy, the achievable reduction fell from 58\% to 14\% and CARMA's advantage over MPC became small (0.6 points). Carbon-aware shifting is physically most valuable when it absorbs low-carbon output that would otherwise be curtailed, for example behind transmission constraints between Scotland and England, when temporal shifts change which plant is marginal, or when sites lie in different markets. Our marginal accounting is a stylized proxy, not an estimate of GB-wide abatement, so it indicates the direction and rough size of this effect rather than its value. In such settings a marginal signal should drive the scheduler. The scheduling framework is unchanged: only the cost coefficients of the flow network change. Idle power works in the same direction. If servers stay powered, deferral saves only their dynamic energy, and carbon-aware scheduling should be paired with right-sizing~\cite{lin2013dynamic}.

\paragraph{Fair internal carbon accounting} Pooling halved the hindsight emissions of the four sites. However, origin-based accounting, which charges each site for its own jobs wherever they run, would have made South Scotland worse off than on its own in every episode. A pool of organizations or business units allocated this way is unstable, because its cleanest members have an incentive to leave. The dual attribution of Theorem~\ref{thm:core} charges each site the carbon shadow price of its demand and credits it for the capacity it contributes. It is budget-balanced, no coalition would prefer to secede, and it is a by-product of the scheduling problem. It suits internal carbon pricing and the sharing of the benefits of carbon-aware scheduling among business units. It is not a substitute for external GHG Protocol reporting, where negative Scope~2 emissions do not exist~\cite{ghg2015scope2}. The players must also be entities that could credibly schedule on their own, such as separate operators, business units with their own capacity, or tenants with reserved shares. Because credits reward contributed capacity, capacity declarations should be audited.

\paragraph{Deployment} CARMA needs a carbon-intensity feed (average or marginal), a demand profile computed from the operator's own job logs and a network-flow solver. It decides within tens of milliseconds and returns a plan that is optimal for its model, which makes its decisions auditable. Absolute savings scale with the fleet: for the 1{,}700-server fleet simulated here, CARMA saves about 16~tCO$_2$ per year relative to myopic control and about 165~tCO$_2$ relative to carbon-agnostic operation, in attributional terms.

\subsection{Limitations and future research}

The study has several limitations. First, the synthetic workloads use a profile learned from the same generator, which favors anticipation. The trace experiment partly addresses this concern, but it covers only four evaluation days of one cluster, arrival times are first task starts, work is based on planned CPU requests, deadlines and origins are synthetic, and jobs are batched. Its 156 episodes vary carbon conditions but share one demand chronology. Second, the regional intensities published by NESO are modeled estimates, used here as realized values, 52 hours of them are missing and filled by persistence, and the marginal analysis relies on a stylized proxy with fixed thresholds, not on a dispatch model or observed curtailment. Third, the facility parameters (PUE, network energy, a 96-core server) are stylized, and embodied emissions are outside the scope except as the floor that makes $\theta$ finite. Fourth, the forecaster excludes weather inputs, so its day-ahead skill is limited. Fifth, bag-of-tasks jobs with free preemption and migration are an idealization: jobs with state, data gravity, bandwidth limits or data-sovereignty constraints would add restrictions that break the pure network structure. Sixth, the guarantees of Theorem~\ref{thm:carma} assume a backstop resource, a positive emission floor and an episode within the planning window, none of which holds in the simulations. The theorem describes the mechanism and the experiments its empirical profile. Its robustness term is shared by all on-time policies and is far from the $\Omega(\sqrt\theta)$ lower bound. Seventh, the core property is exact only for the hindsight benchmark, and the online adjustment is a heuristic. Eighth, Scenario-MPC was tuned and tested with a modest number of scenarios and was not included in the sensitivity analysis because of its cost. Ninth, the additional fleets use production-based lifecycle intensities computed from open generation data, uniform PUE and a capacity layout that mirrors the British one (tested against its reverse), and they cover Europe, the Americas and Oceania only: no open hourly generation data could be retrieved for Asia or Africa, and the New Zealand site is a single national grid.

Several extensions follow naturally. The ghost-demand profile could be replaced by probabilistic arrival forecasts, with the anticipation weight set by a newsvendor-type argument, or by scenario-based planning that treats demand and carbon uncertainty jointly~\cite{bertsimas2004price}. Tightening robustness toward the lower bound while retaining consistency, for example by protecting deferral decisions with reservation prices~\cite{elyaniv2001optimal,lechowicz2025learning}, is an open theoretical question. The framework can also be driven by marginal emission signals, coupled with right-sizing, on-site renewables and storage~\cite{souza2023ecovisor,acun2023carbon} and demand-response programs~\cite{wierman2014opportunities}, and evaluated across several electricity markets.

\section{Conclusions}\label{sec:conclusion}

We studied carbon-aware scheduling of deferrable computing workloads across data-center sites. The offline problem is a minimum-cost flow, which gives exact oracles (after quantization) and fast online planning. Uncertainty about demand alone forces every online algorithm to a competitive ratio of order $\sqrt\theta$, even with perfect carbon forecasts. CARMA, which adds ghost jobs for expected arrivals to the planning network, is deadline-safe under any predictions, optimal under exact predictions and linearly degrading in the prediction error when a backstop exists and episodes fit within the planning window. The dual prices of the pooled problem give a carbon attribution in the core of the hindsight pooling game.

In out-of-sample tests on five grids across four continents, CARMA came within 2.7--9.8\% of the clairvoyant optimum at a load of 0.5, against 4.4--16.3\% for the best baseline without demand information, and within 21.2\% on a trace-derived workload against 27.5\%. A scenario-based stochastic lookahead with the same demand information was level with CARMA or at most 1 point better, but 12--37 times slower. Anticipating demand was worth more than perfect carbon foresight, and correlated forecast errors hurt more than independent ones. Three results qualify these gains. Without migration, reductions are small and near saturation carbon-aware shifting can be worse than carbon-agnostic operation under a service-adjusted metric. Under a stylized marginal-emission proxy, CARMA's advantage over myopic control shrank to 0.6 points for Great Britain. Origin-based and host-based carbon accounting violated coalitional stability in every test week, whereas the dual attribution did so only for the heuristic online adjustment, and then rarely.

Asia and Africa are not covered, because no open hourly generation data could be retrieved. Carbon-aware computing is most effective when schedulers plan for the work that has not yet arrived, not only for the grid conditions that lie ahead, and when its benefits are measured and shared on a sound basis.
